\chapter{Model Framework}

\section{Overall Model Framework}

This thesis develops a dynamic interbank network simulation model to
examine how different loan-matching mechanisms affect the formation of
interbank exposures and the subsequent evolution of systemic risk.
Unlike models that impose a fixed network, the proposed framework
generates creditor--debtor relationships endogenously from banks'
changing balance-sheet conditions and trading decisions.

The simulated banking system consists of one central bank, 20 commercial
banks, and nine shadow banks. The central bank provides the policy-rate
and liquidity environment but does not participate in ordinary interbank
matching. Commercial banks and shadow banks may become lenders,
borrowers, or non-participants according to their liquidity positions,
capital conditions, funding needs, and risk preferences.

Each simulation period represents one business day. The daily process
proceeds in the following order:

\begin{enumerate}
\def\labelenumi{\arabic{enumi}.}
\item
  Liquidity support queued after the preceding day's settlement is
  disbursed, the policy rule reacts to the preceding systemic-risk value,
  and matured central-bank support loans are settled.
\item
  Payments contractually due under all existing interbank contracts are
  combined into one gross liability matrix. Eisenberg--Noe clearing then
  determines actual payments, receipts, and payment shortfalls.
\item
  Contract balances are updated. With rollover disabled, an unpaid
  single-payment balance causes hard default; with rollover enabled,
  unpaid installments become arrears and do not cause default solely
  because of a missed payment. Creditor write-offs and negative-equity
  failures are resolved to a fixed point.
\item
  Regulatory positions are refreshed, and any liquidity support implied
  by the current settlement is queued for the next business day.
\item
  Market conditions, deposit flows, common shocks, and bank-specific rates
  are updated. Active banks are assigned lender, borrower, or
  non-participant roles.
\item
  New interbank loans are formed through either centralized matching or
  decentralized request-for-quote matching.
\item
  Completed trades are registered as dated contracts, funded projects are
  updated, and project returns, defaults, and losses are realized.
\item
  Solvency support is applied when enabled, the post-project
  negative-equity and creditor-write-off cascade is resolved to a fixed
  point, and systemic-risk indicators are recorded.
\end{enumerate}

Newly originated loans do not enter the clearing calculation on the day
on which they are created because their contractual payment dates lie in
future periods. They nevertheless affect the end-of-period exposure
network and the balance-sheet positions carried into the next business
day.

This sequence links bank conditions, loan formation, contract
persistence, payment clearing, and risk measurement within a common
dynamic framework. Matching decisions made in one period create
obligations that may affect liquidity, capital, and default outcomes in
later periods. The detailed bank variables and market updates are
presented in Sections 3.2--3.4, while the two matching procedures are
specified in Chapter~4 and the clearing mechanism is formalized in
Chapter~5.

\section{Bank Types and Node State Variables}

The model distinguishes three types of institutions: a central bank,
commercial banks, and shadow banks. This section defines their functions
and behavioral characteristics within the model. The numerical
composition of the simulated banking system and the type-specific
experimental settings are reported in Section~6.2.

The central bank is treated as a special policy institution. It provides
the policy-rate environment and the interest-rate corridor but does not
participate in ordinary centralized or RFQ matching as a commercial
lender or borrower. It is excluded from the lender--borrower
role-assignment process and from the calculation of the non-central-bank
failure rate. Liquidity support is provided only when the corresponding
policy-support mechanism is enabled.

Commercial banks are the main ordinary participants in the interbank
market. They may become lenders, borrowers, or non-participants
according to their current liquidity positions, capital conditions,
funding needs, investment opportunities, and risk preferences. Their
market roles are therefore determined dynamically rather than fixed by
institutional type.

Shadow banks also participate as ordinary market institutions, but their
funding structure differs from that of commercial banks. They rely more
heavily on interbank and wholesale funding, which may increase their
sensitivity to refinancing and market-liquidity conditions. However,
shadow banks are not assumed to have automatically higher risk appetite
or predetermined network positions. Their counterparties, exposures, and
market roles emerge from their current states and the selected matching
mechanism.

Commercial and shadow banks are subject to the same capital adequacy
threshold, clearing rules, and failure definition in the baseline model.
Institutional heterogeneity therefore arises primarily from
balance-sheet initialization and funding composition rather than from
different regulatory or failure criteria.

At the node level, each bank is described by a set of balance sheet
variables, regulatory indicators, and behavioral parameters. These
variables jointly determine the bank's risk condition, market role,
lending capacity, borrowing demand, and payment ability during clearing.
In the code implementation, these variables are organized into
15-dimensional node features to represent each bank's state in the
current interbank network. The main node state variables are as follows:


\begin{equation}
\begin{aligned}
\mathbf x_{i,t}=(&C_{i,t},\mathrm{LA}_{i,t},\mathrm{CL}_{i,t},\mathrm{IBA}_{i,t},\mathrm{IBL}_{i,t},S_{i,t},a_{i,t},\\
&\mathrm{CAR}_{i,t},\mathrm{LCR}_{i,t},\mathrm{LEV}_{i,t},V_t,r_{i,t}^{L},r_{i,t}^{I},\rho_{i,t},o_{i,t}).
\end{aligned}
\label{eq:node-state}
\end{equation}
\vspace{-0.4\baselineskip}
\noindent where \(C_{i,t}\), \(LA_{i,t}\), \(CL_{i,t}\), \(IBA_{i,t}\), and
\(IBL_{i,t}\) denote the core capital, liquid assets, current
liabilities, interbank assets, and interbank liabilities of bank \(i\) in
period \(t\), respectively. These five variables describe the bank's
balance-sheet position. A bank with more liquid assets and fewer
short-term liabilities is generally better able to meet its payment
obligations, whereas a large interbank liability position creates
additional repayment pressure.

The variables \(S_{i,t}\) and \(a_{i,t}\) describe solvency and operating
status. \(S_{i,t}\) denotes the solvency ratio, while
\(a_{i,t} \in \left\{ 0,1 \right\}\) is the activity indicator. It equals
one when the bank remains active and zero after the bank has failed. An
inactive bank no longer participates in ordinary interbank matching.
Its contracts enter estate resolution; claims against surviving borrowers
may continue as replacement claims held by the central bank.

Capital and liquidity conditions are represented by \(CAR_{i,t}\),
\(LCR_{i,t}\), and \(LEV_{i,t}\). \(CAR_{i,t}\) measures the bank's
capital relative to its risk-weighted assets, and \(LCR_{i,t}\) measures
the liquid resources available to cover expected short-term funding
outflows. \(LEV_{i,t}\) represents the degree of capital backing for the
asset exposure included in its denominator. These indicators influence
risk assessment, lender--borrower role assignment, and the measurement
of systemic risk.

Market conditions and bank behavior are captured by \(V_{t}\),
\(r_{i,t}^{L}\), \(r_{i,t}^{I}\), \(\rho_{i,t}\), and \(o_{i,t}\).
\(V_{t}\) denotes market volatility in period \(t\). The variables
\(r_{i,t}^{L}\) and \(r_{i,t}^{I}\) denote the bank's lending-rate
requirement and project-investment return rate, respectively. Risk
appetite is represented by \(\rho_{i,t}\), while \(o_{i,t}\) denotes the
expected funding outflow rate. Together, these variables affect lending
willingness, borrowing demand, acceptable financing costs, and
vulnerability to liquidity pressure.

The 15-dimensional feature vector links accounting conditions with
market behavior. Liquid assets, current liabilities, and funding
outflows determine the liquidity buffer; capital and leverage measures
describe financial resilience; and interest rates, volatility, and risk
appetite affect transaction decisions. The features are recalculated as
bank conditions change, ensuring that role assignment and risk
evaluation remain connected to the current balance sheet.

Within each business-day cycle, payments due under existing contracts
are cleared before new loans are formed. The resulting cash positions,
together with market conditions, investment returns, funding outflows,
and outstanding contracts, are used to update the node features. New
lending changes interbank assets and liabilities, while payment
shortfalls feed back into liquidity, capital, and activity status. The
node state is therefore dynamic rather than a fixed initial
characteristic.

\section{Market Environment and Cycle Update Mechanism}

A business day is the smallest common time unit used in the simulation.
Policy rates, lending rates, investment returns, contract maturities,
and balance-sheet growth are therefore expressed at daily frequency.
This is not meant to reproduce intraday payment systems; it provides a
consistent order for settlement, matching, and risk measurement and
prevents an annual rate or maturity assumption from being applied at the
wrong scale.

The market environment is an important external factor affecting bank
behavior and systemic risk. In the model, the market environment is
divided into two states: a bull market and a bear market. A bull market
represents a stable environment, where banks have stronger capital and
liquidity conditions, investment returns are more favorable, and market
volatility is lower. A bear market represents a stressed environment,
where capital and liquidity buffers decline, liability pressure
increases, investment returns deteriorate, and market volatility becomes
higher. Distinguishing between these two market states makes it possible
to analyze systemic risk dynamics under both normal and stressed
conditions.

During the initialization stage, the market environment is randomly set
as either a bull market or a bear market. If the market is in a bull
state, the base interest rate is set at a lower level. If the market is
in a bear state, the base interest rate is set at a higher level. The
long-term interest rate is then determined by adding a
market-state-specific spread to the base interest rate. In the code
implementation, bull and bear markets use different base interest rates,
long-term interest rate spreads, loan rate spreads, and investment
return ranges. The market state affects not only external return
conditions, but also interbank borrowing costs and bank investment
behavior.

The market environment is not fixed throughout the simulation but
evolves over time. Each market state remains in effect for a specified
duration. Once that duration has been reached, the next state is
selected randomly between bull and bear conditions, and the duration
counter is reset. This mechanism captures changes in market conditions
across multiple business days without assuming that the financial system
remains in a single state. Market-state transitions subsequently affect
interest rates, investment returns, liquid assets, and liability
pressure.

After payments due under existing contracts have been settled and the
contract book has been updated, the model revises policy rates and
market conditions for the current business day. In a bull market, the
base interest rate fluctuates around a lower level, the long-term
interest rate is calculated by adding the bull-market spread, market
volatility remains low, and liquid assets may receive positive market
adjustments. In a bear market, the base interest rate may move toward a
higher level, the long-term interest rate is calculated by adding the
bear-market spread, market volatility is higher, and liquid assets may
be negatively adjusted. This setting reflects the increase in funding
costs, decline in asset values, and rise in volatility during stressed
market conditions.

For each bank, the market environment affects several state variables.
The market volatility variable is updated according to the current
market state. The bank's loan interest rate is calculated by adding the
corresponding market-state loan spread to the base interest rate. In a
bear market, the loan spread is usually higher, reflecting higher
funding costs. The bank's investment return rate is determined by adding
an investment return spread to the long-term interest rate. Finally,
liquid assets are adjusted according to market conditions, while current
liabilities grow according to the daily-frequency liability growth rate.

Market conditions also affect banks' funding outflow rates. In a bull
market, the outflow rate of non-central banks remains in a lower range.
In a bear market, the outflow rate is set at a higher level. This means
that banks face greater short-term cash outflow pressure during stressed
periods. If a bank's liquid assets are insufficient relative to its
current liabilities and outflow rate, the simulation regards the bank as
facing short-term liquidity pressure and reduces its risk appetite. The
decline in risk appetite further affects the bank's investment and
lending behavior, allowing market stress to feed back into the interbank
network formation process.

Market-state changes are only one source of variation; the simulation
also includes a common liquidity-and-credit shock. On each business day,
the shock occurs with probability 0.035. A shock does not impose an
unbacked cash-only loss. Instead, it represents an equity-neutral deposit
run: liquid assets and current liabilities fall by the same withdrawal
amount. It simultaneously multiplies the same-day funding-outflow rate by
1.25, subject to heterogeneous bank sensitivity and an upper bound of
0.95. The project-default-probability multiplier rises to 2.0 and remains
in force for five business days. This design creates immediate liquidity
pressure and persistent credit deterioration without mechanically causing
synchronous insolvency in both matching regimes.

The framework also considers the influence of project investment on bank
states. Banks may also hold project investment assets while
participating in the interbank market. Project returns are affected by
the market environment. In a bull market, investment returns are more
favorable, while in a bear market, returns are lower and more volatile.
If a bank has a higher risk appetite and expected project returns exceed
funding costs, it may generate opportunity-driven borrowing demand.
Conversely, when market pressure increases or investment returns
decline, banks' willingness to invest and borrow is reduced. Interbank
borrowing demand is driven not only by liquidity gaps, but also by
investment opportunities and market return conditions.

During the periodic update, current liabilities increase according to
the daily liability growth rate, while liquid assets are adjusted
according to market conditions. This setting allows bank balance sheets
to change continuously over time, rather than only when shocks occur. As
liabilities increase period by period, banks need to maintain higher
liquidity buffers. As liquid assets fluctuate with market conditions,
some banks may change from lenders to borrowers. Market environment
changes affect bank role assignment through balance sheet updates.

The market environment and cycle update mechanism also affect the
structure of the interbank network. In a bull market, banks generally
have stronger liquidity positions and higher investment returns, so
liquidity supply and investment-driven borrowing demand may coexist. In
a bear market, liquidity pressure increases, available lending supply
decreases, and borrowing demand is more likely to arise from liquidity
gaps. In the decentralized RFQ version, borrowing demand is further
decomposed into gap demand and expansion demand. Gap demand is
determined by the difference between target liquidity and actual
liquidity, while expansion demand depends on investment returns, loan
rates, risk appetite, and the market environment.

The market environment and cycle update mechanism have three main
functions in this model. They provide dynamic external conditions for
bank balance sheets and market interest rates, allowing bank states to
change over time. They affect lending capacity and borrowing demand by
changing liquid assets, current liabilities, funding outflow rates,
investment returns, and risk preferences. They influence the dynamic
evolution of systemic risk by changing the conditions for interbank loan
formation and clearing payment capacity. The market environment is not
merely a background setting, but an important driving factor linking
individual bank states, loan matching mechanisms, and systemic risk
propagation.

\section{Regulatory and Risk Indicators}

The simulation calculates both bank-level regulatory indicators and
system-level risk indicators in each simulation period. Bank-level
indicators describe the capital position, liquidity condition, solvency,
and activity status of each bank. These indicators also affect role
assignment and trading behavior in the interbank market. System-level
risk indicators measure the overall risk level of the banking system and
serve as the main output variables for the subsequent experimental
analysis.

\subsection{Capital Adequacy Ratio}

The capital adequacy ratio measures a bank's ability to cover risky
assets with capital. In this model, the capital adequacy ratio is the
core indicator used to identify whether a bank is under capital stress.
Let \(\mathbf{C}_{\mathbf{i},\mathbf{t}}\) denote the core capital of
bank \(\mathbf{i}\) in period \(\mathbf{t}\), and let
\(\mathbf{RW}\mathbf{A}_{\mathbf{i},\mathbf{t}}\) denote its
risk-weighted assets. The capital adequacy ratio is defined as:


\begin{equation}
\mathrm{CAR}_{i,t}=\begin{cases}
 \min\left\{1.5,\,\dfrac{\max\{0,C_{i,t}\}}{\mathrm{RWA}_{i,t}+\varepsilon}\right\},
 &\mathrm{RWA}_{i,t}>0,\\[3pt]
 0,&\mathrm{RWA}_{i,t}=0.
 \end{cases}\label{eq:car}
\end{equation}


In this model, risk-weighted assets consist of interbank assets and
project investment assets, subject to a liability-based numerical floor:


\begin{equation}
\mathrm{RWA}_{i,t}=\max\left\{
0.5\,\mathrm{IBA}_{i,t}+\mathrm{PA}_{i,t},
\ 0.08\,\mathrm{TL}_{i,t}
\right\}\label{eq:rwa}.
\end{equation}
\vspace{-0.4\baselineskip}
\noindent where \(\mathbf{IB}\mathbf{A}_{\mathbf{i},\mathbf{t}}\) denotes
interbank assets and \(\mathbf{P}\mathbf{A}_{\mathbf{i},\mathbf{t}}\)
denotes project investment assets. Interbank assets are assigned a risk
weight of 0.5, while project investment assets receive a risk weight of
1.0. Here \(\mathrm{TL}_{i,t}\) is total accounting liabilities, comprising
current, interbank, and termed-out liabilities. The liability-based floor
prevents a bank with a temporarily small measured risky-asset book from
receiving an unrealistically large CAR. The lower bound \(\varepsilon\) prevents the capital adequacy ratio
from becoming undefined when measured risk-weighted assets are zero or
extremely small. In the code implementation, \_safe\_car\_value() also
limits extreme CAR values to maintain numerical stability.

The model defines the capital adequacy ratio threshold as
\(\mathbf{\theta}\). A bank is considered to be under capital stress
when its capital adequacy ratio falls below this threshold:


\begin{equation}
\mathrm{CAR}_{i,t}<\theta.
\end{equation}


In the baseline setting, the threshold is:


\begin{equation}
\theta=0.08.
\end{equation}


This threshold is used not only to identify banks under capital stress,
but also to calculate the capital-stress share and the capital gap ratio
in the systemic risk indicators.

\subsection{Liquidity Coverage Ratio}

The liquidity coverage ratio measures the extent to which a bank's
short-term liquid assets can cover stressed funding outflows. Let
\(\mathbf{L}\mathbf{A}_{\mathbf{i},\mathbf{t}}\) denote the liquid
assets of bank \(\mathbf{i}\) in period \(\mathbf{t}\),
\(\mathbf{C}\mathbf{L}_{\mathbf{i},\mathbf{t}}\) denote current
liabilities, and \(\mathbf{o}_{\mathbf{i},\mathbf{t}}\) denote the
funding outflow rate. The stressed outflow is defined as:


\begin{equation}
\mathrm{Outflow}_{i,t}
 =\bigl(\max\{0,\mathrm{CL}_{i,t}\}+\max\{0,\mathrm{IBL}_{i,t}\}\bigr)
 \max\{o_{i,t},0.10\}\label{eq:outflow}.
\end{equation}


Here \(\mathrm{IBL}_{i,t}\) denotes outstanding interbank liabilities.
 The effective outflow rate has a floor of 0.10, and termed-out liabilities
 are excluded from this short-term liquidity denominator. With
 \(\varepsilon=10^{-9}\), the liquidity coverage ratio is


\begin{equation}
\mathrm{LCR}_{i,t}=\begin{cases}
 \min\left\{3,\,\dfrac{\max\{0,\mathrm{LA}_{i,t}\}}{\mathrm{Outflow}_{i,t}}\right\},
 &\mathrm{Outflow}_{i,t}>\varepsilon,\\[3pt]
 3,&\mathrm{Outflow}_{i,t}\leq\varepsilon,\ \mathrm{LA}_{i,t}>\varepsilon,\\
 0,&\text{otherwise}.
 \end{cases}\label{eq:lcr}
\end{equation}


In this model, the liquidity coverage ratio directly affects a bank's
funding demand and lending capacity. When a bank's liquid assets are
lower than its target liquidity buffer, it is more likely to become a
borrower. When a bank's liquid assets exceed the target buffer, the
excess liquidity can be used as loanable funds. Therefore,
\(\mathbf{LC}\mathbf{R}_{\mathbf{i},\mathbf{t}}\) is not only a
descriptive risk indicator, but also an important variable in bank role
assignment and loan matching.

The function \texttt{\_safe\_lcr\_value()} implements these rules in both
 mechanisms. The initial exported LCR is evaluated at a construction-time
 outflow rate of 0.4 for non-central banks and 0.2 for the central bank.
 Regulatory refreshes, including the refresh before the first EN settlement,
 then use each bank's current outflow rate. Appendix~\ref{app:initial-bank-data}
 distinguishes this initial cache from subsequent state-dependent values.

\subsection{Solvency Ratio and Leverage Ratio}

In addition to the capital adequacy ratio and the liquidity coverage
ratio, the model also uses the solvency ratio and leverage ratio to
describe bank risk conditions.

The solvency ratio measures accounting equity relative to total
 liabilities. Let \(E_{i,t}=\mathrm{LA}_{i,t}+\mathrm{IBA}_{i,t}
 +\mathrm{PA}_{i,t}-\mathrm{TL}_{i,t}\) denote accounting equity before
 truncation, where total liabilities include current, interbank, and
 termed-out liabilities. The solvency ratio is


\begin{equation}
S_{i,t}=\frac{E_{i,t}}{\mathrm{TL}_{i,t}+\varepsilon}\label{eq:solvency}.
\end{equation}
\vspace{-0.4\baselineskip}
\noindent Regulatory core capital is \(C_{i,t}=\max\{0,E_{i,t}\}\)
 while a bank remains active. A lower solvency ratio indicates a smaller
 equity buffer relative to liabilities; a negative value signals negative
 accounting equity. Liquid assets already enter \(E_{i,t}\) and are not
 added to equity a second time.

The model-specific leverage measure relates a bank's core capital to the
liquid and interbank assets included in the denominator. It is defined
as:


\begin{equation}
\mathrm{LEV}_{i,t}=\frac{C_{i,t}}{\max\left\{\varepsilon,\mathrm{LA}_{i,t}+\mathrm{IBA}_{i,t}\right\}}\label{eq:leverage}.
\end{equation}
\vspace{-0.4\baselineskip}
\noindent where \({LA}_{i,t}\) denotes liquid assets and \({IBA}_{i,t}\) denotes
interbank assets. Because capital appears in the numerator, a higher
value indicates stronger capital backing for these asset exposures. This
measure should not be interpreted as a conventional assets-to-equity
leverage multiple, for which a higher value would indicate greater
leverage.

In the simulation, the solvency ratio, capital adequacy ratio, liquidity
coverage ratio, and leverage ratio are updated in each period and
written into bank state variables.

\subsection{Bank Failure and Capital Stress}

The simulation distinguishes bank inactivity from capital stress. Bank
inactivity is represented by \(a_{i,t}\), where \(a_{i,t} = 1\) indicates
that bank \(i\) remains active in period \(t\), while
\(a_{i,t} = 0\) indicates that it has met the failure condition specified
in the daily simulation process. Once a bank becomes inactive, it is
excluded from subsequent ordinary interbank matching. Default resolution
distributes recoveries and closes its original contracts. Claims held by
the failed bank against surviving borrowers are transferred to the central
bank, retaining the surviving borrowers' repayment obligations.

Capital stress does not necessarily imply immediate failure. An active
bank is classified as being under capital stress when
\({CAR}_{i,t} < \theta\). Such a bank may continue participating in the
interbank market, but its capacity to absorb additional losses is
limited. This distinction allows the analysis to record both realized
failures and vulnerabilities that emerge before a bank becomes inactive.

\subsection{Systemic Risk Indicators}

Let \(\mathcal{B}\) denote the set of non-central banks and let
\(\mathcal{B}_{t}^{active}\) denote the subset that remains active in
period \(t\). The function \(\mathbf{1}( \cdot )\) is an indicator
function that equals one when the stated condition is satisfied and zero
otherwise.

The bank failure rate measures the proportion of non-central banks that
have become inactive:


\begin{equation}
\mathrm{FR}_t=\frac{1}{\lvert\mathcal B\rvert}\sum_{i\in\mathcal B}\mathbf{1}\left\{a_{i,t}=0\right\}\label{eq:failure-rate}.
\end{equation}


Let \(\mathcal A_t=\{i\in\mathcal B:a_{i,t}=1\}\) be the active
non-central banks. The capital-breach share is the fraction of active
banks in the recovery-hysteresis low-CAR state:


\begin{equation}
\mathrm{CBS}_t=
\frac{\displaystyle\sum_{i\in\mathcal A_t}h_{i,t}(\theta)}
{\max\left\{1,\lvert\mathcal A_t\rvert\right\}}
\label{eq:capital-stress-share}.
\end{equation}

Here \(h_{i,t}(\theta)=1\) when current CAR is below \(\theta\), or when
the preceding-period CAR was below \(\theta\) and current CAR remains
below \(\theta+\delta\), with \(\delta=0.0005\). Otherwise it is zero.
Failed banks are excluded from both numerator and denominator and are
recorded exclusively by \(\mathrm{FR}_t\). This definition is used by the
executable scorer, Appendix~A, and all figure labels.

For each non-central bank, the minimum capital requirement is defined
as:


\begin{equation}
K_{i,t}^{\mathrm{req}}=\begin{cases}
 \theta\,\mathrm{RWA}_{i,t},&a_{i,t}=1,\\
 K_i^{\mathrm{fail}},&a_{i,t}=0,
 \end{cases}\label{eq:required-capital}
\end{equation}


The failure snapshot \(K_i^{\mathrm{fail}}\) is recorded using the
 policy threshold and regulatory RWA immediately before the default-resolution
 operation closes the bank's interbank book, and remains fixed thereafter.
 The available capital used in CGR is
 \(C_{i,t}^{\mathrm{avail}}=\mathbf{1}\{a_{i,t}=1\}\max\{0,C_{i,t}\}\).
 The corresponding capital gap is


\begin{equation}
\mathrm{Gap}_{i,t}=\max\left\{0,\,K_{i,t}^{\mathrm{req}}-C_{i,t}^{\mathrm{avail}}\right\}\label{eq:capital-gap}.
\end{equation}


The system-level capital gap ratio is:


\begin{equation}
\mathrm{CGR}_t=\frac{\displaystyle\sum_{i\in\mathcal B}\mathrm{Gap}_{i,t}}{\varepsilon+\displaystyle\sum_{i\in\mathcal B}K_{i,t}^{\mathrm{req}}}\label{eq:capital-gap-ratio}.
\end{equation}
\vspace{-0.4\baselineskip}
\noindent Whereas
\(\mathrm{CBS}_t\) records how widely capital stress is distributed,
\(\mathrm{CGR}_t\) measures the aggregate severity of the resulting capital
shortfalls.

The composite systemic risk indicator is defined as:


\begin{equation}
\mathrm{SR}_t=w_1\mathrm{FR}_t+w_2\mathrm{CBS}_t+w_3\mathrm{CGR}_t\label{eq:systemic-risk}.
\end{equation}


The baseline weights are:


\begin{equation}
\left(w_1,w_2,w_3\right)=\left(0.5,0.3,0.2\right).
\end{equation}



\begin{equation}
w_1+w_2+w_3=1.
\end{equation}


The resulting value of \(\mathrm{SR}_t\) is restricted to the interval
\(\left\lbrack 0,1 \right\rbrack\). The three components are also
retained separately so that changes in the composite measure can be
attributed to realized failures, the breadth of capital stress, or the
severity of capital shortfalls.

\subsection{Role of Indicators in the Simulation}

The regulatory indicators connect bank balance sheets with market
behavior and systemic risk measurement. At the bank level, \(S_{i,t}\),
\({CAR}_{i,t}\), \({LCR}_{i,t}\), \({LEV}_{i,t}\), and \(a_{i,t}\) form
part of the node state used in role assignment and post-clearing
updates. Liquidity conditions affect borrowing demand and lending
capacity, while capital and solvency conditions determine how strongly a
bank can withstand investment losses and payment shortfalls. At the
system level, \(\mathrm{FR}_t\), \(\mathrm{CBS}_t\), and \(\mathrm{CGR}_t\) distinguish
realized failure from emerging and accumulated capital pressure. Their
weighted combination, \(\mathrm{SR}_t\), provides the main measure used to
compare centralized matching and decentralized RFQ.

\section{Period Definition and Dynamic Process}

Each simulation period represents one business day. Market conditions,
interest rates, investment returns, balance-sheet variables, contract
maturities, and systemic-risk indicators are therefore expressed at a
common daily frequency.

Conceptually, each period connects lagged policy support, settlement of
inherited obligations, market and shock updates, formation of new
interbank loans, project outcomes, solvency support, recursive default
resolution, and risk measurement. The state produced at the end of one
period becomes the initial state of the next period.

Outstanding contracts connect these stages across time. Payments that
become due in period \(t\) enter the current clearing process, whereas
principal and interest scheduled for later periods remain in the
contract book. New loans formed in period \(t\) enter the end-of-period
exposure network but do not enter the due-payment liability matrix until
their scheduled payment dates.

The conceptual transition can therefore be summarized as:

previous-period state \(\rightarrow\) queued support and policy
\(\rightarrow\) one EN settlement \(\rightarrow\) default fixed point
\(\rightarrow\) market and shock update \(\rightarrow\) daily matching
\(\rightarrow\) project outcomes and solvency support \(\rightarrow\)
post-project default fixed point \(\rightarrow\) risk measurement.

The executable simulation algorithm, including initialization, daily
iteration, Monte Carlo aggregation, and early stopping, is specified in
Section~6.4.
